% ============================================================
% OP_Q11_pagelevel_verification_v2.tex  (FINAL, for GPT sign-off)
% Page-level verification of the Q11 Euclidean spinor/OS
% conventions against the ORIGINAL constructions, read in full:
%   [L] Luscher, Commun. Math. Phys. 54, 283-292 (1977) (PDF, 10 pp.)
%   [OS] Osterwalder-Schrader, Helv. Phys. Acta 46, 277-302 (1973)
%        (e-periodica scan, 27 pp., DOI 10.5169/seals-114484)
% plus the previously verified [KU] arXiv:1012.5601 (full text),
% [KU'] arXiv:1005.3751 (excerpt), hep-lat/0306014 App. B (excerpt),
% Williams CMP 38, 65 (1974) (introduction).
% NO preprint edits made; proposed edits listed in Sec. 4 for
% approval.
% ============================================================
\section*{Q11 page-level convention verification (final)}

\subsection*{1. Findings, item by item (with page/equation cites)}

\textbf{(1) Hermitian Euclidean gammas} (ours: eq.~q11\_eucl\_gamma).
[L] footnote~5, p.~286: ``The $\gamma$-matrices used are euclidean
ones, i.e. $\{\gamma_\mu,\gamma_\nu\}=2\delta_{\mu\nu}$;
$\gamma_\mu^+=\gamma_\mu$'' --- \emph{verbatim} our convention;
explicit chiral representation in his eq.~(32).
[OS] p.~282: $\gamma_E^0=\gamma^0$, $\gamma_E^j=+i\gamma^j$
(Bjorken--Drell side), hermitian, Euclidean Clifford algebra;
$p_E=\sum p_l\gamma_l^E$.
Note: [OS] Wick-map runs from the $(+,-,-,-)$ BD convention while
our eq.~q11\_wick\_map runs from $(-,+,+,+)$; both land on the same
hermitian Euclidean algebra. Our conventions block already labels
the map as a convention-matching statement only --- no change
needed.

\textbf{(2) Fermionic time reflection with $\gamma_0$ and species
exchange} (ours: eq.~q11\_eucl\_theta\_spinor).
[OS] eqs.~(4.1)--(4.3), pp.~290--291: unitary involution $\Theta$
(``time inversion followed by charge conjugation'') exchanging the
doubled species $\Psi^{(1)}\!\leftrightarrow\!\Psi^{(2)}$ with the
matrix $\gamma_0$ and mode conjugation.
[L] p.~284, below eq.~(4): $\theta$ = euclidean time reflection,
conjugation $\mathcal O\to\mathcal O^+$ with the explicit example
$\psi^+=\gamma_0\bar\psi$ (Grassmann conjugation per Berezin, his
ref.~[8]).
[KU] eqs.~(16)--(17): verbatim our form.
Verdict: our $\theta\psi=\bar\psi\gamma_0$,
$\theta\bar\psi=\gamma_0\psi$ is the Berezin-variable form of
exactly this structure.

\textbf{(3) Antilinearity and order reversal} (ours:
eq.~q11\_eucl\_theta\_spinor, third relation).
[L] eq.~(5): $(\mathcal O_1,\mathcal O_2)
=\langle\theta(\mathcal O_1^+)\mathcal O_2\rangle$, antilinear in
the first slot; Grassmann conjugation is an anti-involution
(order-reversing), per Berezin.
[OS] antilinearity realised at form level via
$(\,\Theta X,Y\,)$ with mode conjugation in eqs.~(4.2)--(4.3).
[KU] eqs.~(18)--(19): verbatim.

\textbf{(4) Form order, weight, pairing} (ours:
eq.~q11\_graded\_os\_form with $e^{-S_E/S_0}$, pairing $+S_E$).
[L] eq.~(4): $\langle\theta(\mathcal O^+)\mathcal O\rangle\ge0$,
reflected element \emph{first}; weight $e^{+A}$ (his eqs.~(1),(6):
expectations $Z^{-1}\!\int\!\mathcal DU\,\mathcal D\bar\psi\,
\mathcal D\psi\,\phi_1\cdots\phi_n\,e^{A}$).
[OS] \S4.1: $\langle X,Y\rangle=(\Theta X,Y)$, reflected element
first; positivity by Lemma~4.1 (p.~292), extended to interacting
Green's functions in \S6 (announced p.~278).
[KU]~(1012.5601): $\langle\theta(F)F\rangle\ge0$ with $e^{-S}$;
[KU$'$]~(1005.3751): $e^{+A}$ --- both action signs occur even
within one author pair.
Verdict: the original [L] uses the combination
$(e^{+A},\ \Theta\text{-first})$; ours is
$(e^{-S_E},\ F\text{-first})$ --- exactly two global signs apart
(pairing convention $\times$ quadratic-level reordering), i.e. the
same projector-positive branch.
The appendix~(b) statement (``Standard formulations state
reflection positivity as $\langle(\Theta F)F\rangle\ge0$, with
action-sign, measure-ordering, and pairing conventions that vary
between sources'') is verified \emph{verbatim} at page level.

\textbf{(5) Grassmann positivity in the originals; commuting case.}
[L] Propositions~1--2 (pp.~287--289): strictly positive transfer
matrix and OS positivity for Grassmann Wilson fermions
($0<K<1/6$).
[OS] Lemma~4.1 + \S6: positive semi-definite form, positivity
axiom verified.
Neither source considers the commuting/symmetric alternative ---
recorded as absent; the converse (negative-mode) side of Q11.1a is
\emph{our} added content, as claimed.

\textbf{(6) Doubling.}
[OS] \S3.2, pp.~284--285, verbatim: the naive single-field ansatz
fails (``non-hermitian, hence a fortiori not positive definite''),
``A way out is to double the number of fields'' ---
$\Psi^{(1)},\Psi^{(2)}$ independent.
[L] independent classical Grassmann $\psi,\bar\psi$; operator
$\psi_n$ without canonical anticommutator (his eq.~(12)).
Our independent Berezin pair $(\psi,\bar\psi)$ is this doubling.

\textbf{(7) Non-hermitian Euclidean action $\leftrightarrow$
$\Theta$-adjoint.}
[OS] Introduction item~III (p.~278): the fermionic Euclidean
action is non-hermitian and ``the adjoint transformation
$V\to V^*$ is related to Euclidean time inversion'' --- the same
structure as our conventions block~(iii) (Euclidean adjoint defined
through $\Theta$).

\subsection*{2. Conclusion}
All checklist items (L1)--(L5), (O1)--(O5) are verified with
page/equation citations; no mismatches found.
The trial conventions of \S qs\_reflection\_positivity are a
faithful Berezin-variable presentation of the original
Osterwalder--Schrader / L\"uscher constructions, with the order
and action-sign related by the documented two-global-sign
equivalence.
The residual review condition (c) of Appendix~app:q11\_os\_test is
discharged.

\subsection*{3. Proposed status change (for GPT sign-off)}
Q11.1a: candidate Level~A $\to$ \textbf{Level~A} (mathematical
result for the free minimal Euclidean Dirac kernel; ECT application
Q11.1b unchanged at Level~B conditional; Q11.2/Q11.3 unchanged;
Q11.4 open).

\subsection*{4. Edit sites upon approval (no edits made yet)}
\begin{enumerate}
  \item Appendix (c): rewrite as completed verification with the
    page/equation citations above (sources: Luscher1977,
    OsterwalderSchraderFermi1973, KikukawaUsui2011).
  \item Proposition Q11.1a status line: remove ``candidate'' and
    the review-condition clause; state verification completed
    (Appendix).
  \item First mention in the OP-Q11 decomposition paragraph:
    ``candidate Level~A ... pending'' $\to$ ``Level~A mathematical
    result (conventions verified against the original OS/L\"uscher
    constructions; Appendix)''.
  \item Status longtable row Q11.1a: ``A candidate (math.)'' $\to$
    ``A (math.)''; comment updated.
  \item fermion\_open item (ii): ``candidate mathematical result''
    $\to$ ``Level~A mathematical result (free kernel)''.
  \item Part III graph node q11\_spinstat: ``Q11.1a cand.~A'' $\to$
    ``Q11.1a A''; regenerate figures from
    fig\_partIII\_vertical\_v2.gv.
  \item Companion: one-word alignment (``convention-fixed candidate
    calculation'' $\to$ ``convention-fixed calculation'' +
    ``verified against the original Euclidean Fermi-field
    constructions'').
  \item Consistency sweep: any remaining ``candidate'' tied
    specifically to Q11.1a (grep); Q11.1b/AN-statuses untouched.
\end{enumerate}

\subsection*{5. Honest limits}
The verification covers the defining data and the positivity
statements of the originals; it does not re-derive their theorems.
[L]/[OS] prove the Grassmann-side positivity for their kernels
(Wilson lattice / free continuum with doubling); our Q11.1a
statement is for the free minimal continuum kernel in our
conventions, with our own two-family calculation and our own
converse; the originals are used to verify \emph{conventions},
not to import theorems.
